
\section{Notebook 18: Spectral Memory and Low-Rank Operator Correction}

Notebook 18 studies whether the two-step transition residual from the normalized
prime-gap state process is low-rank.  Let
\[
z_n = \frac{p_{n+1}-p_n}{\log p_n}
\]
be the normalized prime gap, and let $s_n$ be a quantile-discretized state of
$z_n$.  The first-order transition operator is
\[
P_{ij}=\Pr(s_{n+1}=j\mid s_n=i),
\]
and the empirical two-step operator is
\[
P^{(2)}_{ik}=\Pr(s_{n+2}=k\mid s_n=i).
\]
The Markov prediction is $\widehat{P}^{(2)}=P^2$, so the two-step memory
residual is
\[
M=P^{(2)}-P^2.
\]
We decompose this residual as
\[
M=U\Sigma V^\top
\]
and define a rank-$k$ correction
\[
P^{(2)}_k=P^2+U_k\Sigma_kV_k^\top.
\]

For this run, the Markov $L_2$ error is
\[
0.00793828,
\]
while the rank-4 corrected $L_2$ error is
\[
0.00124468.
\]
The rank-4 correction reduces error by
\[
0.843205.
\]
The leading singular mode accounts for energy share
\[
0.557366,
\]
and the rank needed to capture 90\% of residual energy is
\[
3.
\]
A shuffle baseline gives mean top singular value
\[
0.0113954,
\]
compared with the real top singular value
\[
0.0474118.
\]
These diagnostics test whether higher-order transition memory is structured as a
small number of spectral modes rather than spread uniformly across transitions.
